\section{Introduction}
\label{sec:intro}

In 2021, the United States maternal mortality rate reached 32.9 deaths per 100,000
live births, the highest level since 1965~\citep{Toy2023WSJ,Harris2023}. Between 2018 and 2022,
Non-Hispanic Black women died at 2.8 times the rate of White women, and
disorders related to pregnancy, including hemorrhage and venous
complications, were the leading underlying cause of death, accounting for
17.4\% of 6{,}283 pregnancy-related deaths~\citep{Chen2025}. These figures
highlight the role of the coagulation system in maternal outcomes. Blood coagulation is a complex enzymatic cascade in which a small
tissue factor (TF) stimulus triggers a series of serine protease activation reactions
that ultimately produce thrombin, the central effector of
hemostasis~\citep{Mann2011,Butenas2009}. Thrombin cleaves fibrinogen into fibrin,
activates platelets, and amplifies its own production through positive feedback via
Factors~V, VIII, and XI, while simultaneously initiating its own shutdown through the
thrombomodulin--protein~C anticoagulant pathway~\citep{Hockin2002,Bravo2012}. Pregnancy
alters this hemostatic balance, shifting toward a hypercoagulable state
that prepares for delivery~\citep{Hellgren2003,Grouzi2022,Bremme2003}.
Procoagulant factors including fibrinogen, Factor~VIII, and von Willebrand factor
increase substantially across gestation, while natural anticoagulants such as
antithrombin and protein~S decline. Two conditions further shift this balance in ways that are not fully understood.
Polycystic ovary syndrome (PCOS), a hormonal and metabolic disorder affecting
approximately 6.6\% of reproductive-age women~\citep{Azziz2004}, has been associated with
impaired fibrinolysis and a prothrombotic tendency~\citep{Palomba2015}. Preeclampsia (PE), a hypertensive
disorder that complicates 2--8\% of pregnancies worldwide~\citep{Abalos2014}, is
characterized by endothelial dysfunction, platelet activation, and altered thrombin
generation~\citep{Dusse2011}. Mathematical models of the coagulation cascade have provided mechanistic insight
into patient-specific thrombin generation~\citep{Luan2007,Luan2010}. These range
from the original Hockin--Mann stoichiometric framework~\citep{Hockin2002} through
extensions incorporating the thrombomodulin--protein~C pathway~\citep{Bravo2012}
to the Brummel-Ziedins 2012 (BZ2012) model with detailed activated protein~C
(APC)-mediated Factor~Va inactivation
kinetics~\citep{BrummelZiedins2012}. However, applying these models to pregnancy
cohorts requires longitudinal datasets capturing coagulation factors, inhibitors,
and thrombin generation assay (TGA) measurements across gestation, data that is
rare and expensive to collect.

The core barrier is sample size. Longitudinal pregnancy coagulation studies
typically enroll tens of patients with complete follow-up across multiple trimesters,
yielding datasets where the number of features~$p$ (coagulation factors, TGA
parameters, viscoelastic measurements) exceeds the number of patients~$n$. Rare
complications compound this problem. Small longitudinal cohorts inevitably contain
only a handful of patients with any given condition, far too few for
condition-specific statistical analysis. This $n < p$ regime limits conventional
approaches because covariance matrices are rank-deficient, cross-validation is
unreliable, and overfitting is nearly inevitable~\citep{Hastie2009}. Synthetic data generation addresses this by augmenting small real datasets with statistically faithful synthetic
records, enabling downstream analyses (mechanistic modeling, clinical comparisons,
hypothesis generation for rare complications) that would otherwise be
infeasible~\citep{Gonzales2023,Kuhnel2024}. \rtwo{An original motivation for this work was
to support deep learning models in both directions: predicting TGA responses from coagulation
factors and estimating factor profiles compatible with an observed TGA response.} However, existing methods face
limitations in the small-cohort longitudinal setting. Sampling from a fitted
multivariate normal (MVN) distribution is singular when $n < p$ and must be
regularized~\citep{Ledoit2004}, introducing bias that distorts the joint distribution.
MVN also assumes Gaussian marginals and linear correlations, and provides no mechanism
for conditional generation: it cannot amplify a specific subpopulation while preserving
its signatures. More expressive models such as generative adversarial networks (GANs) and
variational autoencoders (VAEs)~\citep{Goodfellow2014,Kingma2014,Xu2019CTGAN}, including recent longitudinal
extensions~\citep{Kuhnel2024}, require substantially larger training sets and are
prone to mode collapse at small $n$~\citep{DeMitri2025} (we confirm this empirically for Conditional Tabular GAN (CTGAN) at
$n{=}23$ in this work). A generative method is needed that can operate
directly on the geometry of a small dataset without \rone{estimating a full
covariance model}.

Stochastic Attention (SA) is such a framework. \rone{Rather than fitting a full
covariance model or training a neural generator,} SA treats the stored patient
profiles themselves as the model. Building on modern Hopfield network
theory~\citep{Ramsauer2021}, each profile becomes a memory pattern in a
continuous energy landscape\rone{. This energy induces a continuous probability
distribution that is exactly a finite Gaussian mixture in the reduced space,
with one nonzero-variance component centered on each stored pattern. Sampling
from this distribution therefore produces continuous variations around the
stored patterns rather than exact copies}. Because sampling occurs in a
principal component analysis (PCA)-reduced linear subspace anchored on the observed cohort, SA never forms a
full $p \times p$ covariance matrix and so avoids both the rank deficiency that
limits MVN and the mode collapse that limits GAN/VAE approaches when $n < p$.
We recently introduced a multiplicity-weighted extension of the Hopfield energy
that attaches a per-pattern weight to each stored profile~\citep{Varner2026Multiplicity},
which provides exactly the lever that MVN lacks: continuous interpolation
between unconditioned generation and targeted amplification of a rare
subpopulation at inference time, without retraining. What remains open, and what
this study addresses, is whether SA can faithfully generate \emph{longitudinal}
trajectories from a small cohort. In earlier proceedings work, we paired our
mechanistic coagulation pipeline with per-visit MVN sampling to produce
synthetic populations~\citep{Varner2023JuliaCon}; that approach captured
single-visit marginals but, by fitting each visit independently, generated
patients whose trajectories across gestation were statistically incoherent. A
generator that respects the joint structure across visits is what a small
longitudinal pregnancy cohort actually demands.

In this study, we used multiplicity-weighted SA to generate complete
longitudinal patient profiles from a small pregnancy coagulation cohort
($K{=}23$ patients, 72 features per visit, 3 visits) that includes rare
clinical subgroups (PCOS, $n{=}3$; Developed~PE, $n{=}5$). Our pipeline
concatenates each patient's multi-visit profile into a single vector, applies
PCA for dimensionality reduction, and generates new patients via
multiplicity-weighted Langevin dynamics on the Hopfield energy landscape, with
a direction-magnitude decomposition that preserves the natural variance
structure of continuous clinical measurements. We then asked whether the
synthetic patients reproduced marginal feature distributions, the cross-visit
covariance structure, condition-specific signatures of rare clinical
subgroups, and the biological consistency obtained when the profiles were processed by
\rboth{a separately specified, generator-blind ordinary differential equation (ODE) model
of the coagulation cascade that was calibrated on the same real
cohort}~\citep{BrummelZiedins2012}. As a downstream-utility test, we further
asked whether a mechanistic model calibrated entirely on synthetic patients
could predict held-out real patient outcomes. \rboth{As a proof-of-concept,} the
results show that SA recovers the statistical structure of the real cohort and
matches its mechanistic plausibility at this sample size. If this generalizes, the implication is that the bottleneck
for studying rare obstetric and pediatric conditions may be shifting from
cohort size toward cohort fidelity: a few dozen carefully phenotyped
longitudinal patients, augmented through SA, may be enough to support
mechanistic and statistical analyses that have traditionally relied on much
larger cohorts.
